% General Introduction. Sources: the authors' account of ARC 3 and ARC 4 (user, 2026-10-26,
% including their written feedback to the organisers), Chapter 2 (specification, Table 2.1).
\chapter*{General Introduction}
\addcontentsline{toc}{chapter}{General Introduction}
\markboth{General Introduction}{}

\section*{Context}
Student robotics competitions have become, for many electronics students, the place where
the subjects taught separately at university first meet in a single object: a robot
combines a microcontroller, power electronics, sensors, wireless communication, mechanical
design and software, and it must work on a given day, on a track it has not seen before. In
Algeria, the Algerian Robot Cup (ARC) is one of these competitions. Each edition opens with
a qualification stage, in which the robot must follow a line through a maze, followed by a
main event whose task changes from one year to the next and is described in a technical
specification published by the organisers.

\section*{Problem statement}
The authors took part in two consecutive editions of the ARC, ARC~3 and ARC~4, and in
several other competitions, without success. The fourth edition illustrates the
difficulties they met. Initially planned for 14~October~2025, it was held on 5~December~2025.
The robot was in good working order at the beginning of October; in the weeks that
followed, parts were added and removed in anticipation of the requirements, and a week
before the event the ESP32-S3 controlling the robot was damaged and had to be replaced.
The test day that was to precede the competition was cancelled, so the robot first met the
real track during the competition itself. On the day, the gripper broke, the arm's yaw servo
stopped working, the base of the arm cracked and then broke, one motor driver
behaved erratically and the other was disconnected when the robot was closed. Several
elements of the track also differed from the technical specification, which had been read
and prepared in detail beforehand.

Looking back on both editions, the authors trace most of these difficulties to the same
pattern. For each competition, the robot was largely rebuilt from scratch: chassis, sensor
mounts and wiring, and in the case of ARC, two robots instead of one. The time this took
left little for testing, and the uncertainty about the task led to a robot more
complicated than needed and modified until the last day, with parts that had never been
tested together. A competition robot built this way is new every time, and so are its
faults.

This raises the question this work sets out to answer: \emph{can a single, ready-to-use
mobile robot base, with a standard way of mounting and wiring its sensors, let a team
reconfigure its robot for each competition instead of rebuilding it?}

\section*{Objectives}
The answer proposed in this thesis is BLOC-BOT, a modular mobile robot platform. Its
objectives are:
\begin{itemize}
  \item a differential-drive robot based on an ESP32-S3 microcontroller, with wheel
        encoders, an ultrasonic distance sensor and an inertial sensor, built from
        components that can be bought locally at low cost;
  \item a handheld controller, based on an ESP32, that drives the robot over the ESP-NOW
        wireless protocol without any computer or phone, and displays the data the robot
        sends back;
  \item a robot that stops on its own when it no longer receives commands;
  \item a mechanical standard, a 3D-printed plate with a regular hole pattern and printed
        holders called \emph{blocs}, so that sensors and modules can be added, moved or
        exchanged without redesigning the chassis.
\end{itemize}
The full list of requirements, and the state of each one at the time of writing, is given
in Table~\ref{tab:requirements}.

\section*{Organisation of the thesis}
The thesis is organised in three chapters. Chapter~\ref{ch:generalities} gives the
background on which the work builds: mobile robots and their locomotion, educational,
competition and modular robots, microcontrollers, and wireless links between them.
Chapter~\ref{ch:design} presents the study and design of the system: the specification,
the overall architecture, and the choice and dimensioning of each component of the robot,
of the controller and of the wireless link. Chapter~\ref{ch:realisation} describes the
realisation: the wiring, the firmware of both boards, the problems met during development
and how they were solved, the tests carried out and the cost of the platform. A general
conclusion summarises the results and the limits of the work and proposes directions for
future work.
